\documentclass[10pt,a4paper]{amsart}
%---------------------------------------------------------%
% FAVOR NÃO ALTERAR O PREÂMBULO
%---------------------------------------------------------%

\usepackage{amsmath,amssymb,amsfonts,amsthm,amscd}
\usepackage{enumerate}
\usepackage{lipsum}
\usepackage[english]{babel}

%------------------------------------------%
% FAVOR NÃO ADICIONAR NENHUM PACOTE %
%------------------------------------------%


\begin{document}

\title{\textsf{Boolean algebras of infinite graphs and Tangle Spaces}}
\author{\textbf{Violeta Martins de Freitas}, \underline{Leandro Aurichi}} 

\maketitle
  
\vspace{-0.7cm}
{\small

% Este arquivo deve ter uma página, contando as referências

I will present the recently submitted paper I co-authored with Gustavo Boska on a Boolean algebra based approach to studying infinite graphs. In it, we define two algebras which are able to capture the connectivity structure of a given infinite graph. These two algebras, when applied the famous Stone duality functor on, give rise to a number of topological spaces that have been studied in the literature of infinite graphs for the last 60 years. This provides a new approach to defining and working with these structures. In the paper, we also pursue attempts to functorialize these constructions. I will also present how the ideas developed through the course of this paper led to the work I am developing right now on tangle theory. The abstract theory of tangles was developed by graph theorists in the last 10 years to a wide success of different applications in combinatorics. I will show how tangles are intimately connected to ultrafilters and how that can lead to a host of new and interesting questions for tangles in set theory and logic.

}

\textit{Acknowledgments:} We would like to thank CAPES for Financial Support. \\

\begin{thebibliography}{99} 

% Colocar referências no seguinte modelo 

%\newline{}[1] 
\bibitem{key1} Violeta Mar, Gustavo Boska, \emph{Describing ends and tangles (and their edge variants) through Boolean algebras and functors}, \textit{pre-print on arXiv:2606.17201}, \textbf{2026}.
    
%\newline{}[2]
\bibitem{key2} Jay Kneip \emph{Tangles and where to find them}, \textit{Dissertation}, \textbf{2020}. 
    

\end{thebibliography}


{\scriptsize

Violeta Mar {ICMC-USP}:
\email{violetamar@usp.br}\\
Leandro Aurichi {ICMC-USP}:
\email{aurichi@icmc.usp.br}
} 

\end{document}